\chapter{Protocolo experimental}
\label{cap:protocolo-experimental}

El protocolo distingue entre el diseño aprobado y la evidencia ya ejecutada. La especificación autoritativa de la campaña final es \path{experiments/final_experiment_campaign.json}; los artefactos históricos se conservan para contexto, pero no satisfacen una celda final por coincidencia aproximada. El perfil científico \texttt{main-v1} está congelado y la campaña no modifica arquitectura, recompensa, contrato canónico ni hiperparámetros MAIN.

\section{Preguntas experimentales}

La campaña responde a siete preguntas ordenadas:

\begin{enumerate}[label=\textbf{P\arabic*.},leftmargin=*]
    \item ¿reproduce una MAIN fresca la capacidad de aprendizaje intradistribución del pipeline actual?;
    \item ¿coinciden la evaluación persistida y una validación directa e independiente del modelo?;
    \item ¿qué parte del resultado aleatorio puede asociarse a duplicados y cuánto cambia al separar días?;
    \item ¿cómo varía el comportamiento al aumentar el tamaño de entrenamiento y su presupuesto proporcional?;
    \item ¿cuánto cambia al variar solo la semilla del modelo sobre una partición fija?;
    \item ¿qué sensibilidad aparece al retener cuatro escenarios dirigidos?;
    \item ¿los patrones observados son específicos de QRDQN o también aparecen en Random Forest, y cómo se transfieren a la captura de laboratorio?
\end{enumerate}

La Figura~\ref{fig:campana-final} muestra las dependencias entre bloques. MAIN es a la vez experimento primario y fuente de los controles auxiliares; los alias reutilizan una ejecución física solo cuando la especificación declara identidad exacta.

\begin{figure}[htbp]
\centering
% Figura conceptual: matriz y dependencias, sin resultados medidos.
% Fuente: experiments/final_experiment_campaign.json.
\begin{tikzpicture}[node distance=5mm and 5mm]
  \node[cajaResaltada, text width=3.3cm] (main)
    {\textbf{MAIN fresca}\\full, random\\$3\,000\,000$ pasos\\s42 / m42};

  \node[cajaProceso, text width=3.3cm, right=of main] (day)
    {\textbf{Día completo}\\Lun--Mié $\to$ Jue--Vie\\$3\,000\,000$ pasos};
  \node[cajaProceso, text width=3.3cm, right=of day] (holdouts)
    {\textbf{4 holdouts}\\Web, Infiltration,\\PortScan, DDoS\\$1\,000\,000$ pasos};

  \node[cajaProceso, text width=3.3cm, below=8mm of main] (ladder)
    {\textbf{Escalera}\\100k, 250k, 500k,\\1M, 2M, full};
  \node[cajaProceso, text width=3.3cm, right=of ladder] (seeds)
    {\textbf{Semillas}\\1M filas; $1\,324\,741$ pasos\\model 42--46};
  \node[cajaModelo, text width=3.3cm, right=of seeds] (rf)
    {\textbf{RF emparejado}\\full, 1M, día\\y 4 holdouts};

  \node[cajaDato, text width=11.3cm, below=8mm of seeds] (aux)
    {\textbf{Auxiliares de MAIN}\\validación directa $\cdot$ bootstrap 10k $\cdot$ duplicados $\cdot$ etiquetas barajadas 10k $\cdot$ Fase~2 fresca};

  \draw[flechaAzul] (main) -- (day);
  \draw[flechaAzul] (day) -- (holdouts);
  \draw[flechaAzul] (main) -- (ladder);
  \draw[flechaAzul] (ladder) -- (seeds);
  \draw[flechaAzul] (holdouts) -- (rf);
  \draw[flechaAzul] (day) -- (rf);
  \draw[flechaAzul] (main) -- (aux);

  \node[etiquetaTenue, text width=11.3cm, below=2mm of aux]
    {22 ejecuciones primarias físicas $\rightarrow$ 24 puntos lógicos mediante 2 alias; 5 trabajos auxiliares.};
\end{tikzpicture}

\caption[Visión general de la campaña experimental final]{Relación entre la MAIN fresca, la partición completa por días, la escalera de tamaños, la sensibilidad a semilla, los cuatro escenarios retenidos, las comparaciones RF y las validaciones auxiliares. Las flechas indican reutilización de artefactos o comparación emparejada, no resultados. Fuente: elaboración propia a partir de \path{experiments/final_experiment_campaign.json}.}
\label{fig:campana-final}
\end{figure}

\section{Datos, unidad de análisis y fases}

La unidad es una fila de flujo. La Fase~1 emplea los ocho CSV de CICIDS2017 después de limpieza, mapeo binario y transformación al esquema canónico. La Fase~2 aplica el modelo final de forma offline a una captura etiquetada del laboratorio del autor. En ningún caso se actúa sobre una red ni se retroalimenta el entorno con la decisión.

La composición por día de CICIDS2017 se mantiene como contexto del protocolo porque explica por qué una separación temporal modifica simultáneamente prevalencia y repertorio de ataques. La Figura~\ref{fig:composicion-dia} procede de conteos comprometidos, no de la futura campaña.

\begin{figure}[htbp]
\centering
\includegraphics[width=0.94\textwidth]{f6_composicion_dia.pdf}
\caption[Composición por día de CICIDS2017]{Filas y proporción de tráfico benigno/ataque por CSV de CICIDS2017. Fuente: conteos deterministas conservados en \texttt{memoria/figuras/data\_composicion\_dia.json}.}
\label{fig:composicion-dia}
\end{figure}

\section{Preprocesamiento y control de fuga}

El pipeline construye primero una representación canónica \emph{sin escalar}. La caché aprobada almacena esta representación y metadatos de procedencia para evitar repetir la carga y coerción de los CSV. No almacena valores normalizados: el escalador depende de la partición y debe ajustarse dentro de cada ejecución.

Para una ejecución dada, el orden es:

\begin{enumerate}[leftmargin=*]
    \item cargar o validar la caché canónica sin escalar;
    \item seleccionar filas y construir la partición usando \texttt{split\_seed};
    \item ajustar \texttt{StandardScaler} exclusivamente sobre entrenamiento;
    \item transformar entrenamiento y prueba con ese estado;
    \item entrenar con \texttt{model\_seed};
    \item persistir modelo, escalador, percentiles, métricas y procedencia.
\end{enumerate}

La separación entre semillas evita que un estudio de estabilidad del modelo cambie inadvertidamente el test. Los identificadores de flujo, IP, tiempo y puertos específicos excluidos por el contrato canónico no se reintroducen desde la caché. Los controles de manifiesto y firma impiden reutilizar una caché creada con otro conjunto de fuentes o versión de esquema.

\section{Perfil congelado \texttt{main-v1}}

El perfil MAIN fija la representación de 152 dimensiones, la arquitectura \texttt{[1024,1024,512]}, 200 cuantiles, tasa de aprendizaje $5\times10^{-5}$, $\gamma=0$, 20 pasos de gradiente por ciclo y la recompensa \texttt{tp=1.5}, \texttt{fp=-2.0}, \texttt{fn=-5.0}, \texttt{omission=0.0}. La ejecución completa usa $3\,000\,000$ de pasos. Estos valores proceden del perfil congelado, cuya firma se registra en la especificación; no se vuelven a ajustar durante la campaña.

La MAIN final \texttt{qrdqn\_main\_random\_full\_s42\_m42} usa el conjunto completo, partición aleatoria estratificada, \texttt{split\_seed=42}, \texttt{model\_seed=42} y $3\,000\,000$ de pasos. Es una ejecución fresca: la MAIN histórica de 2026 no se promueve ni se copia para ocupar esta celda.

\section{Matriz primaria}

\subsection{Partición aleatoria y día completo}

La MAIN usa una división aleatoria estratificada 80/20. Este diseño conserva la prevalencia aproximada y ofrece el punto intradistribución común para QRDQN y RF. Su objetivo es verificar capacidad de aprendizaje, no generalización externa.

El contraste temporal \texttt{qrdqn\_day\_full\_s42\_m42} entrena con lunes, martes y miércoles y prueba con jueves y viernes. Usa el perfil completo y $3\,000\,000$ de pasos. Sustituye como evidencia temporal principal a la Comprobación~C histórica, que era una red proxy de $30\,000$ pasos.

\subsection{Escalera de tamaños}

La escalera mantiene la partición aleatoria y fija ambas semillas en 42. Incluye $100\,000$, $250\,000$, $500\,000$, $1\,000\,000$, $2\,000\,000$ de filas de entrenamiento y el conjunto completo. Los presupuestos son, respectivamente, $132\,474$, $331\,185$, $662\,370$, $1\,324\,741$, $2\,649\,482$ y $3\,000\,000$ de pasos. El punto completo es un alias de MAIN y no requiere una segunda ejecución física.

La curva permitirá describir eficiencia dentro de este protocolo. Como tamaño y pasos cambian juntos, no se interpretará como efecto causal puro del número de filas.

\subsection{Sensibilidad a la semilla del modelo}

El estudio fija $1\,000\,000$ de filas, $1\,324\,741$ pasos y \texttt{split\_seed=42}. Solo cambia \texttt{model\_seed} entre 42, 43, 44, 45 y 46. El punto 42 es alias del peldaño de un millón. La salida final reportará la distribución y los puntos individuales, no solo una media.

\subsection{Cuatro escenarios retenidos dirigidos}

El estudio retiene exactamente cuatro CSV: Web Attacks, Infiltration, PortScan y DDoS. Cada QRDQN se entrena durante $1\,000\,000$ de pasos sobre los restantes CSV, con semillas de partición y modelo fijadas en 42. El propósito es cubrir escenarios de naturaleza distinta con un coste experimental controlado. No se presenta como leave-one-CSV-out exhaustivo de los ocho ficheros.

\section{Random Forest emparejado}

Random Forest utiliza la misma representación canónica, filas de entrenamiento, conjunto de prueba y métricas de cada contraste. La campaña incluye RF para la MAIN aleatoria completa, el punto aleatorio de un millón, la partición completa por días y los cuatro escenarios retenidos. No recibe presupuesto de pasos, pero sí un contrato de partición y semillas explícito.

Esta línea base permite separar una dificultad del dominio de un efecto específico de QRDQN. Las comparaciones se realizarán por pares; no se mezclan con puntuaciones publicadas bajo otros filtros o particiones.

\section{Controles auxiliares}

La MAIN fresca alimenta cinco trabajos auxiliares:

\begin{itemize}[leftmargin=*]
    \item \textbf{validación directa}: recarga modelo y preprocesamiento, reproduce la partición y verifica predicciones y conteos;
    \item \textbf{bootstrap}: $10\,000$ remuestreos estratificados con semilla 12345 sobre el test fijo;
    \item \textbf{duplicados}: cuantifica duplicados internos y coincidencias exactas entre entrenamiento y prueba;
    \item \textbf{etiquetas barajadas}: control breve de $10\,000$ pasos que debe perder la señal predictiva cuando se destruye la relación con la etiqueta;
    \item \textbf{Fase 2 fresca}: aplica MAIN, escalador y percentiles finales a la captura de laboratorio.
\end{itemize}

El \emph{bootstrap} no reentrena el modelo y no mide varianza entre semillas. El análisis de duplicados tampoco atribuye por sí mismo causalidad al rendimiento. Ambos se interpretan junto con los contrastes que cambian la distribución.

\section{Jerarquía de evaluación}

La Figura~\ref{fig:escalera-validacion} organiza las afirmaciones posibles. El nivel no aumenta porque una métrica sea más alta, sino porque la separación entre entrenamiento y prueba es más exigente. La Fase~2 añade una fuente distinta, aunque su captura siga siendo cerrada y generada por el autor.

\begin{figure}[htbp]
\centering
\resizebox{\textwidth}{!}{% Fuentes: runs/validation/VAL_checks_A_20260212_235443/validation_results.json
% (Comprobacion A: accuracy=0.9939, n_samples=10000);
% runs/validation/VAL_checks_B_20260212_235736/validation_results.json
% (Comprobacion B: shuffled_accuracy=0.4773, baseline_accuracy=0.5227,
% leakage_detected=false); runs/validation/duplicate_analysis_seed42.json
% (cross_split_leakage_feature_only.pct_of_test=24.8624);
% runs/validation/VAL_checks_C_20260213_004847/validation_results.json
% (Comprobacion C: recall_attack=0.5295374374439701, train Mon/Tue/Wed,
% test Thu/Fri); runs/cicids2017/baseline_random_forest_comparison/
% rf_cicids2017_canonical_20260628_024735__day_split/metrics.json
% (recall_attack=0.08135289329152054); runs/phase2/
% P2v2_pred_20260610_161231_MAIN/metrics.json (block_rate=0.252364,
% accuracy=0.991862).
\begin{tikzpicture}[every node/.style={text width=3.6cm}]

% --- Primer tramo (menos exigente): 1 -> 2 -> 3 ---
\node[caja] (s1) at (0,0) {\textbf{1. Comprobación A}\\evaluación directa};
\node[etiquetaTenue, below=1mm of s1] (s1v) {Exactitud 0.9939\\(10\,000 muestras)};

\node[caja] (s2) at (4.0,0) {\textbf{2. Comprobación B}\\etiquetas barajadas};
\node[etiquetaTenue, below=1mm of s2] (s2v) {Exactitud 0.4773 vs.\ 0.5227 base\\(sin fuga detectada)};

\node[cajaPendiente] (s3) at (8.0,0) {\textbf{3. Duplicados}\\prueba vs.\ entrenamiento};
\node[etiquetaTenue, below=1mm of s3] (s3v) {24.86\,\% de la prueba\\ya vista en entrenamiento};

% --- Segundo tramo (más exigente), tras un descanso: 4 -> 5 -> 6 ---
\node[caja] (s4) at (0,5.4) {\textbf{4. Comprobación C}\\partición por día};
\node[etiquetaTenue, below=1mm of s4] (s4v) {Recall de ataque 0.5295\\(red proxy, no pesos MAIN)};

\node[caja] (s5) at (4.0,5.4) {\textbf{5. Línea base RF}\\mismo protocolo};
\node[etiquetaTenue, below=1mm of s5] (s5v) {Recall de ataque 0.0814\\(partición por día)};

\node[cajaResaltada] (s6) at (8.0,5.4) {\textbf{6. Fase 2}\\tráfico de laboratorio};
\node[etiquetaTenue, below=1mm of s6] (s6v) {Bloqueo 25.24\,\%; exactitud 99.19\,\%\\(validez externa limitada)};

\draw[flechaAzul] (s1) -- (s2);
\draw[flechaAzul] (s2) -- (s3);
\draw[flechaAzul] (s3.north) -- ++(0,0.7) -| (s4);
\node[etiquetaTenue, text width=1.8cm] at (4.0,2.3) {(descanso)};
\draw[flechaAzul] (s4) -- (s5);
\draw[flechaAzul] (s5) -- (s6);

\node[etiqueta, text width=3.6cm, below=5mm of s1v] {menos exigente};
\node[etiqueta, text width=3.6cm, above=5mm of s6] {más exigente};

\end{tikzpicture}
}
\caption[Jerarquía de evaluación y generalización]{Jerarquía de evidencia: control intradistribución, separación por día, escenarios retenidos y transferencia a laboratorio. La sensibilidad a tamaño y semilla caracteriza estabilidad, pero no constituye por sí misma un nivel de validez externa. Fuente: elaboración propia.}
\label{fig:escalera-validacion}
\end{figure}

\section{Métricas y reglas de lectura}

Se reportan matriz de confusión, exactitud, exactitud balanceada, MCC, precisión, \emph{recall} y F1 por clase, además de tasas de falsos positivos y falsos negativos. La clase ATTACK es positiva. Para seguridad, \emph{recall} de ataque y tasa de falsos negativos reflejan ataques omitidos; para impacto, la tasa de falsos positivos refleja benignos bloqueados. La exactitud nunca se interpreta aislada.

Las curvas de entrenamiento describen optimización, no evaluación. Los intervalos se presentan con su unidad de incertidumbre. Las comparaciones por escenario incluyen tamaño y prevalencia del test para evitar que una media oculte distribuciones distintas.

\section{Fase 2: inferencia offline y diagnósticos}

La inferencia fresca carga el modelo, escalador, percentiles y nombres de características de MAIN. El mapeo construye la observación canónica, registra qué columnas fuente están disponibles, sanea valores no finitos de acuerdo con el contrato y aplica el escalador persistido. No vuelve a ajustar ningún estado con la captura de laboratorio.

Los diagnósticos incluyen presencia de características, valores no finitos, valores estandarizados extremos, recorte respecto a percentiles de entrenamiento y tasa de bloqueo. Si existen etiquetas válidas se calculan las mismas métricas binarias. La Figura~\ref{fig:pipeline-fase2} refleja que la fuente del modelo será la MAIN final y que las salidas mantienen procedencia.

\begin{figure}[htbp]
\centering
% Fuente: scripts/predict_real_traffic_v2.py (orden de etapas, ficheros de
% artefactos, recortes, predicción determinista y diagnósticos);
% src/canonical_schema.py (dimensión canónica 152 = 76 características + 76 de máscara de presencia).
\begin{tikzpicture}[
    node distance=5mm and 5mm,
    every node/.style={text width=5cm},
  ]

  % --- Columna izquierda ---------------------------------------------
  \node[cajaDato] (pcap) {Captura pcap\\{\footnotesize\sffamily\color{tesisGris}laboratorio doméstico cerrado}};
  \node[cajaProceso, below=of pcap] (flujos) {Extracción de flujos\\{\scriptsize\ttfamily\color{tesisGrisMedio}flows.csv}};
  \node[cajaProceso, below=of flujos] (canon) {Mapeo a esquema canónico (152 dim.)\\{\footnotesize\sffamily\color{tesisGris}76 características + 76 de máscara de presencia}\\{\scriptsize\ttfamily\color{tesisGrisMedio}map\_to\_canonical()}};
  \node[cajaProceso, below=of canon] (perc) {Recorte por percentiles del train (p0.5/p99.5)\\{\scriptsize\ttfamily\color{tesisGrisMedio}train\_percentiles.npz}};

  % --- Columna derecha -------------------------------------------------
  \node[cajaProceso, right=20mm of pcap] (scaler) {Escalado con \texttt{StandardScaler} del MAIN\\{\footnotesize\sffamily\color{tesisGris}reutilizado, NO reajustado}\\{\scriptsize\ttfamily\color{tesisGrisMedio}scaler.joblib}};
  \node[cajaPendiente, below=of scaler] (clipz) {Recorte z opcional\\{\footnotesize\sffamily\color{tesisGris}\texttt{-{}-clip-z}, si se indica}};
  \node[cajaModelo, below=of clipz] (qrdqn) {QRDQN determinista\\{\footnotesize\sffamily\color{tesisGris}model.predict(deterministic=True)}};
  \node[cajaResaltada, below=of qrdqn] (out) {Predicciones permitir / bloquear\\{\footnotesize\sffamily\color{white}+ métricas y diagnósticos z}\\{\scriptsize\ttfamily\color{white}predictions.csv, metrics.json}};

  % --- Flechas de flujo --------------------------------------------------
  \draw[flecha] (pcap) -- (flujos);
  \draw[flecha] (flujos) -- (canon);
  \draw[flecha] (canon) -- (perc);
  \draw[flechaAzul] (perc.east) -- ++(6mm,0) |- (scaler.west);
  \draw[flecha] (scaler) -- (clipz);
  \draw[flecha] (clipz) -- (qrdqn);
  \draw[flecha] (qrdqn) -- (out);

  % --- Nota inferior -------------------------------------------------
  \node[etiquetaTenue, below=6mm of out, text width=10.5cm] (nota)
    {Inferencia puramente offline: el script solo escribe \texttt{predictions.csv}, \texttt{metrics.json} y \texttt{config.json}; no existe bloqueo activo del tráfico.};

\end{tikzpicture}

\caption[Pipeline de inferencia offline de la Fase 2]{La MAIN final proporciona modelo, escalador, percentiles y contrato; la captura de laboratorio proporciona flujos y, cuando procede, etiquetas. La salida incluye predicciones, métricas y diagnósticos, sin actuar sobre la red. Fuente: elaboración propia a partir del pipeline actual de Fase~2.}
\label{fig:pipeline-fase2}
\end{figure}

\section{Contrato de artefactos y cierre de campaña}

Cada ejecución primaria del esquema~3.0 persiste, como mínimo, \texttt{config.json}, \texttt{metrics.json}, \texttt{environment.json}, \texttt{run\_summary.json}, \texttt{timings.json}, monitorización, modelo, escalador, percentiles, nombres de características, manifiesto y sumas SHA-256. Los trabajos auxiliares usan contratos equivalentes adaptados a su salida. El \texttt{RUN\_ID} identifica el directorio y enlaza configuración, evidencia y agregado.

La campaña ejecuta 22 trabajos primarios físicos que representan 24 puntos lógicos mediante dos alias, más cinco trabajos auxiliares. El cierre exige que todos estén completos, sean coherentes con \texttt{main-v1} y pasen la agregación. Los alias no duplican modelos ni exportaciones.

Los artefactos sellados se copian al destino de evidencia junto con un archivo comprimido y su suma. Git conserva especificación, índices, métricas y manifiestos adecuados para revisión; Git LFS se reserva para binarios rastreados que exceden el flujo normal. Los datos originales, cachés grandes, directorios de ejecución activos y exportaciones voluminosas no se incorporan por defecto al historial.

\section{Plataforma experimental y preflight}

La metodología final utiliza lenguaje neutral: \emph{plataforma GPU experimental}. Antes de ejecutar se comprueban commit, limpieza del árbol, disponibilidad de datos y LFS, firma del perfil, caché, espacio, memoria, CUDA, importaciones, capacidad de crear artefactos y un ensayo corto no científico. Esta verificación determina si la plataforma puede ejecutar el protocolo; no produce resultados de tesis.

El hardware, sistema operativo, controlador, CUDA, GPU, CPU, RAM y paquetes finales solo se describirán desde los \texttt{environment.json} de las ejecuciones aceptadas. No se anticipa un proveedor ni una GPU concreta. Los detalles de la plataforma histórica permanecen únicamente en el Anexo~\ref{cap:anexo-entorno} como procedencia de evidencia previa.

\section{Estado del protocolo}

El diseño, la implementación y los contratos de artefactos están preparados. La ejecución científica permanece pendiente. Esta separación permite redactar el método como texto casi final y, al mismo tiempo, impide que una orden preparada, un \emph{dry-run} o un artefacto histórico se conviertan en evidencia final.

\evidencenote{La campaña no se ha ejecutado durante esta revisión de la memoria. El Capítulo de Resultados solo podrá poblarse después de validar el conjunto completo de artefactos contra la especificación congelada.}
